\begin{table}[htbp]
\centering
\caption{TN1 -- tỉ lệ đạt J tốt nhất (\%) theo số hạt và số vòng trên Taiwan Default, c~=~3}
\label{tab:e5-luoi-taiwan_kiem_chung}
\footnotesize
\setlength{\tabcolsep}{3pt}
\begin{tabular}{@{}lrrrrrrrrr@{}}
\toprule
& \multicolumn{3}{c}{10 hạt} & \multicolumn{3}{c}{20 hạt} & \multicolumn{3}{c}{30 hạt} \\
\cmidrule(lr){2-4} \cmidrule(lr){5-7} \cmidrule(lr){8-10}
Phiên bản & T = 20 & T = 30 & T = 50 & T = 20 & T = 30 & T = 50 & T = 20 & T = 30 & T = 50 \\
\midrule
PSO gốc (1995) & 100 & 100 & 100 & 100 & 100 & 100 & 100 & 100 & 100 \\
PSO-W (1998) & 100 & 100 & 100 & 100 & 100 & 100 & 100 & 100 & 100 \\
PSO-S (2022) & 100 & 100 & 100 & 100 & 100 & 100 & 100 & 100 & 100 \\
PSO-LOG (2022) & 100 & 100 & 100 & 100 & 100 & 100 & 100 & 100 & 100 \\
QPSO (2004) & 100 & 100 & 100 & 100 & 100 & 100 & 100 & 100 & 100 \\
\multicolumn{10}{@{}l}{\textit{Đối chứng FCM (U ngẫu nhiên): 100\%}} \\
\multicolumn{10}{@{}l}{\textit{Đối chứng FCM (điểm ngẫu nhiên): 100\%}} \\
\multicolumn{10}{@{}l}{\textit{Đối chứng FCM (k-means++): 100\%}} \\
\bottomrule
\end{tabular}
\par\smallskip\parbox{\linewidth}{\scriptsize Ghi chú: 30 seed mỗi ô; T: số vòng PSO trước khi FCM chạy tới hội tụ.}
\end{table}
